%% nota.tex.j2 — template LaTeX da nota informativa da PIM-PF.
%% Renderizado por render_latex.py com delimitadores Jinja customizados (ver o modulo).
\documentclass[11pt,a4paper]{article}
\usepackage[utf8]{inputenc}
\usepackage[T1]{fontenc}
\usepackage{lmodern}
\usepackage[brazil]{babel}
\usepackage[margin=2cm]{geometry}
\usepackage{booktabs}
\usepackage{graphicx}
\usepackage{xcolor}
\usepackage{helvet}
\renewcommand{\familydefault}{\sfdefault}
\usepackage{array}
\setlength{\parindent}{0pt}
\setlength{\parskip}{6pt}
\pagestyle{empty}

\definecolor{azulpim}{HTML}{2166AC}
\definecolor{cinzapim}{HTML}{555555}
\definecolor{vermelhoaviso}{HTML}{B2182B}

\begin{document}


{\color{azulpim}\Large\textbf{Nota PIM-PF — Produção Industrial}}\\[0.2em]
{\color{cinzapim}\small Mês de referência: \textbf{abril 2026}}

\vspace{0.6em}
{\large\textbf{Produção industrial avançou 0,7\% em abril 2026 (m/m-1, com ajuste sazonal)}}\\[0.2em]
{\color{cinzapim}Interanual: +2,7\% · Acumulado no ano: +1,7\% · Em 12 meses: +0,7\%}

\vspace{0.6em}
Em abril 2026, a produção industrial brasileira avançou 0,7\% ante o mês anterior na série com ajuste sazonal, com o índice dessazonalizado em 106,6 pontos (base 2022=100). Na comparação interanual, a indústria avançou 2,7\%; no acumulado do ano, avançou 1,7\%. O avanço mensal reforça a trajetória de alta também observada no acumulado em 12 meses.

\vspace{0.8em}
{\color{azulpim}\textbf{As quatro leituras (indústria geral)}}
\begin{center}
\begin{tabular}{>{\raggedright}p{9.5cm} r}
\toprule
\textbf{Leitura} & \textbf{Variação} \\
\midrule
Variação mensal (m/m-1, com ajuste sazonal) & +0,7\% \\
Variação interanual (m/m-12) & +2,7\% \\
Acumulado no ano & +1,7\% \\
Acumulado em 12 meses & +0,7\% \\
\bottomrule
\end{tabular}
\end{center}

\vspace{0.3em}
{\color{azulpim}\textbf{Por seções}}
\begin{center}
\begin{tabular}{l r r r r}
\toprule
\textbf{Seção} & \textbf{m/m-1} & \textbf{Interanual} & \textbf{Ano} & \textbf{12 m} \\
\midrule
Indústrias extrativas & +3,1\% & +10,6\% & +9,3\% & +7,3\% \\
Indústrias de transformação & +0,3\% & +1,2\% & +0,3\% & -0,5\% \\
\bottomrule
\end{tabular}
\end{center}

\vspace{0.2em}
\noindent\begin{minipage}{0.5\linewidth}\centering
\includegraphics[width=\linewidth]{C:/Users/joaoz/tutorial 5 programação/output/grafico_serie.pdf}
\end{minipage}\hfill
\begin{minipage}{0.48\linewidth}\centering
\includegraphics[width=\linewidth]{C:/Users/joaoz/tutorial 5 programação/output/grafico_categorias.pdf}
\end{minipage}

\vspace{0.4em}
{\color{azulpim}\textbf{Destaques entre atividades (interanual)}}\\
\textbf{Altas:} Fabricação de coque, de produtos derivados do petróleo e de biocombustíveis (+13,3\%); Fabricação de produtos de borracha e de material plástico (+3,8\%); Fabricação de produtos alimentícios (+3,2\%).\\
\textbf{Baixas:} Fabricação de outros equipamentos de transporte, exceto veículos automotores (-7,9\%); Fabricação de máquinas e equipamentos (-7,0\%); Confecção de artigos do vestuário e acessórios (-6,5\%).

\vspace{0.3em}
Entre as grandes categorias econômicas (m/m-1, com ajuste sazonal), bens de capital avançou 0,1\%; bens intermediários avançou 1,5\%; bens de consumo duráveis recuou 3,2\%; bens de consumo semiduráveis e não duráveis recuou 0,2\%.

\vfill
{\color{cinzapim}\footnotesize
\rule{\linewidth}{0.4pt}\\
Fonte: IBGE — Pesquisa Industrial Mensal / Produção Física (PIM-PF Brasil), via SIDRA — tabelas 8888 e 8887.
Extração em 11/06/2026.
}

\end{document}